%    (C) Copyright 2023, 2025 Anthony D. Dutoi
% 
%    This file is part of Qode.
% 
%    Qode is free software: you can redistribute it and/or modify
%    it under the terms of the GNU General Public License as published by
%    the Free Software Foundation, either version 3 of the License, or
%    (at your option) any later version.
% 
%    Qode is distributed in the hope that it will be useful,
%    but WITHOUT ANY WARRANTY; without even the implied warranty of
%    MERCHANTABILITY or FITNESS FOR A PARTICULAR PURPOSE.  See the
%    GNU General Public License for more details.
% 
%    You should have received a copy of the GNU General Public License
%    along with Qode.  If not, see <http://www.gnu.org/licenses/>.
%
\documentclass[12pt,letterpaper]{article}
\setlength{\pdfpagewidth}{\paperwidth}
\setlength{\pdfpageheight}{\paperheight}
\setlength{\textheight}{24cm}
\setlength{\textwidth}{17cm}
\addtolength{\topmargin}{-2.5cm}
\addtolength{\hoffset}{-1.6cm}
\usepackage{xcolor}
\usepackage{listings}
\usepackage{hyperref}

\definecolor{codegreen}{rgb}{0,0.6,0}
\definecolor{codegray}{rgb}{0.5,0.5,0.5}
\definecolor{codepurple}{rgb}{0.58,0,0.82}
\definecolor{backcolour}{rgb}{0.95,0.95,0.92}
\lstdefinestyle{codestyle}{
    basicstyle=\ttfamily\footnotesize,
    breakatwhitespace=false,
    breaklines=true,
    keepspaces=true,
    numbers=left,
    numbersep=5pt,
    showspaces=false,
    showstringspaces=false,
    showtabs=false,
    backgroundcolor=\color{backcolour},
    commentstyle=\color{codegreen},
    keywordstyle=\color{magenta},
    numberstyle=\tiny\color{codegray},
    stringstyle=\color{codepurple},
}
\lstset{style=codestyle}

\begin{document}

\title{TensorNet Subpackage:\\Usage and Coding Documentation}
\author{\copyright~ Copyright 2023, 2025 Anthony D. Dutoi}
\date{}
\maketitle

This file is part of \texttt{Qode}.
 
\texttt{Qode} is free software: you can redistribute it and/or modify
it under the terms of the GNU General Public License as published by
the Free Software Foundation, either version 3 of the License, or
(at your option) any later version.
 
\texttt{Qode} is distributed in the hope that it will be useful,
but WITHOUT ANY WARRANTY; without even the implied warranty of
MERCHANTABILITY or FITNESS FOR A PARTICULAR PURPOSE.  See the
GNU General Public License for more details.
 
You should have received a copy of the GNU General Public License
along with \texttt{Qode}.  If not, see \url{http://www.gnu.org/licenses/}.



		\section{Quick Start}

The following code is imported directly from \texttt{tensornet/tests/tutorial.py}.
It is in the \texttt{tests/} directory because running it and checking the output
 according to the embedded comments serves as a quick (but not exhaustive) test
 that the subpackage is working correctly.
It is named \texttt{tutorial.py} because it is heavily commented to illustrate
 basic usage.

\lstinputlisting[language=Python]{../tests/tutorial.py}

		\section{More Stuff}

In here should also be the outline of \lstinline{tensor_base} and an example of a backend interface class.

Also, we should start with some overview of the file/directory structure.

	\subsection{Original introduction to \texttt{tensors.py}}

We have 4 different types of top-level tensors
  \lstinline{contraction_expression}
  \lstinline{tensor_sum}
  \lstinline{tensor_network}
  \lstinline{primitive_tensor}
All of these are built off of \lstinline{tensor_base}.  The latter three are built off of \lstinline{resolved_tensor} (which
is based on \lstinline{tensor_base}) indicating that they are completely well-formed.  \lstinline{contraction_expression}
behaves like a tensor so long as it is well formed, but it could be incomplete.  The basic problem
is that, when connecting a bunch of tensors with @, it is not possible to automatically discern
when the string is finished (ie, when all intended target free indices are present).  Only the user can decide
when to "close" the string, but we don't want to make the user do this explicitly.  Therefore,
we allow this incomplete class, whose *resolution* is a \lstinline{tensor_network} (it always resolves to a
\lstinline{tensor_network}). This is triggered by taking another action on it that is not @ or *.  A \lstinline{tensor_sum}
can only contain types \lstinline{primitive_tensor} and \lstinline{tensor_network}.  An attempt to include another
\lstinline{tensor_sum} will cause it to be expanded out as a single sum (and \lstinline{contraction_expression} will 
be automatically resolved to \lstinline{tensor_network}).  A \lstinline{tensor_network} can only contain
type \lstinline{primitive_tensor}, and including one network (or \lstinline{contraction_expression}) into another network causes explicit expansion
into a single network.  Taking this together, this means that a \lstinline{tensor_network} cannot contain a 
\lstinline{tensor_sum}.  An attempt to put a \lstinline{tensor_sum} into a \lstinline{tensor_network} will result in the hierarchy
being swapped to a sum of networks.  The reason for this is purely practical.  At the moment,
I cannot think of a way of finding the best path to evaluate a network that contains sums.

	\subsection{Original introduction to \texttt{tensor\_network} in \texttt{tensors.py}}

Barebones theory (written much later after a forensic debug battle).  A \lstinline{tensor_network} object contains
two fundamental pieces of information (and other incidental info).  One is a list of contractions.
Each contraction itself is a list of two-tuples; each two-tuple identifies a tensor and the index
of that tensor involved in the contraction.  So if a contraction contains two two-tuples, then 
two indices are contracted with one another, but there might be more for unusual contractions.
The ordering of the list of contractions is irrelevant, as is the ordering of the list of two-tuples
that defines a given contraction. The other piece of information is a list of free indices of the result,
which correspond to uncontracted indices of the tensors in the network.  This list is ordered, and each
free index is itself a list of two-tuples with the same tensor-index structure as the two-tuples that
define a contraction.  In the most usual case, each such list corresponding to a free index will be of
length one, but if there is more than one tensor index that corresponds to a single free index it is because
those indices are set equal to each other and reduced to a single free index.


	\subsection{Original introduction to \texttt{primitive\_tensor\_factory} in \texttt{\_\_init\_\_.py}}

Philosopy:  A backend (in the sense of the value of the argument backend below) is a class that
contains static methods that perform basic operations on objects that represent tensors of a specific
type. These will henceforth be called backend tensors.  Tensornet then uses the backend functions to
manipulate the backend tensors according to what is encoded by its high-level syntax.  Objects of the
class defined below are factories for the primitive tensornet tensors which each wrap a backend tensor
(and never modifies it).  An important point is that a given backend can have multiple factories
depending on how what a user is doing.

The creation of a tensornet primitive tensor should only take place via a factory (which is why
that class name is "hidden" here with the underscore convention) because this is the mechanism by
which a backend tensor is associated with its backend functions and the utility for contraction creation.
In the most basic case, this creation is done by the .init() method (not \lstinline{.__init__}, which creates the 
factory).  The method .zeros() is a convenience function for creation of a tensornet primitive storing
a new backend tensor (via a backend method) so the user need not create such a zero tensor first and
then "wrap" it (\lstinline{.scalar_tensor()} is similar and was added for completion, but commenting it out essentially
documents that we have not yet found a use case ... the necessity for the backend to be able to create
such a thing however has to do with obeying expectations about type uniformity).

Finally we come to the arguments \lstinline{data_wrap} and \lstinline{data_unwrap}, the method .data(), and the idea of multiple
factories for the same backend.  It is important that tensornet works only with a notion of backend
tensors and the functions it is given for them without knowing anything further about their properties.
But the user may be working with some abstracted backend tensors.  Although they contain, or
somehow represent the existence of, a specific concrete tensor (eg, numpy, tensorly, pytorch), they 
do not have the usual expected behavior (e.g., [] might not even get an element).  We will call such 
things meta tensors, it might be very inconvenient for the user to make their in-house meta tensor behave 
exactly like a concrete tensor, right down to letting a numpy.ndarray (etc) be instantiated from it.  It
will also be very inconvenient if the user has to always use the unbound function raw() and then afterward
extract the concrete data for cooperation with other linear algebra libraries.  The converse is also true
that it will be annoyingly repetitive to always create the meta tensor from a concrete tensor from elsewhere,
and then pass that to the factory.  So, therefore, a factory can be created that takes the functions
\lstinline{data_wrap} and \lstinline{data_unwrap}.  Although their actions are not constrained as such, the idea is that \lstinline{data_wrap}
is a function that takes a concrete tensor and returns an object of the meta tensor type with which
backend works, and \lstinline{data_unwrap} does the opposite, extracting and returning (by reference) the concrete
tensor represented (eg, contained within) the meta tensor.  The convenience function .data() called
on a tensornet tensor causes evaluation of the (potential?) network (sum?) and returns directly the concrete
tensor associated with the result.  By default, \lstinline{data_wrap} and \lstinline{data_unwrap} are just pass-throughs,
which means that .data() is just an alias for the unbound function raw(), which forces evaluation and
returns a backend tensor (which is typically already a concrete tensor type).  Since this automatic
wrapping and unwrapping might not always be the desired behavior, more than one factory can be defined,
per need.

Addendum 1:  There is no requirement that a meta tensor is not itself also a concrete type, and that
this concrete type might be initialized from multiple other concrete types, which might not even be
different than its own type.  So if we interpret the concrete and meta layers to both be tensorly,
where the meta layer can be initialized also from numpy or pytorch (etc), then there it is a perfectly
clean for the user to provide as \lstinline{data_wrap} a function that makes a tensorly tensor from some input
which may or may not be of tensorly type), and to leave \lstinline{data_unwrap} as the default pass-through.

Addendum2:  One question that can be asked about the design is: why am I giving in a separated
backend with functions that perform manipulations on a tensor, rather than just insisting that backend
tensors have certain properties, like += (which would also make the insides of tensornet prettier).
The answer is that one might have enough serious work to do in the abstracted/meta backend already,
without having to define yet more operations that make the thing look almost like a tensor, but not
completely.  One can imagine various inheritance and virtual method schemes, but one is always going
to end up in one of two places: 1. Forcing the user to make a new class for their backend tensors
(and then maybe always, even when using "normal" tensors), or 2. automating the creation of the 
backend tensors with the expected operators in terms of some other/divorced set of instructions.
In case 2, we are back where we started from the user point of view, though the insides of tensornet
look nicer (for the price of another file/class and layer of abstraction).

Addendum 3: The \lstinline{.__init__()} method used to take a \lstinline{copy_data} argument which meant that the factory should
copy the provided primitive before storing it as an argument for further use.  Probably unusual, but
a user might do this if they intended on modifying that primitive for some other use while it was still
in use by the tensornet code (you never know ...).  This can now be done by the user with the \lstinline{data_wrap}
option.




		\section{To Do}

document thing about \lstinline{scalar_tensor}
OR
a new tensor class called \lstinline{scalar_tensor} which is the result of any operation that reduces to
a scalar (recognized from shape).  the difference between this an the usual scalar type is that it is mutable (which 
is very important for evaluating \lstinline{tensor_sum}

implement \lstinline{__str__} for all tensors
check that internal data is copied in the event of trivial network (single \lstinline{primitive_tensor}) and that scalar is 1

\begin{lstlisting}
backend:
	allow variable field type (could have a name?)
        backend base class with more things implemented for "normal" backends
	        scalar_value and scalar_tensor = pass
		raw_check = pass (user can check that input tensors are as expected)
		increment -> +=
		mult -> *
		subscript -> []
		shape -> .shape
		str -> str
		field = float
\end{lstlisting}


clean up documentation. Some good doc sentences...

The string of tensors to be contracted often
consists of more then two, and so it would be incorrect to assume that a complete and
valid tensor has been described after every use of the @ operator.  Therefore, it is not
safe to assume that the use of @ results in something that is a tensor in every way
(eg, it may not even have a valid shape, due to missing free indices).  By the same token,
since the = operator in python only means "bind" (as opposed to, say, C++), there is no
way for the user to indicate when a tensor description is "done".
(correction, now use \lstinline{__call__()})
  The dark corner here
is that the result of several @ operations remembers the names of its contraction indices
and if that result is bound to a variable, it could be @-ed with more tensors, inadvertently
using the same letters, leading to unexpected results.  This is probably not a big deal,
however, because this usage pattern is also unorthodox, since, after binding the result
of an @ string to a name, that object should itself be subjected to the call operation ()
referencing the indices of the finished product (with no memory of labels used for internal
contractions) before being contracted with other tensors.

	\subsection{Questions about using \texttt{opt\_einsum}}

\url{https://joss.theoj.org/papers/10.21105/joss.00753}

Clearly there are advantages to the fact that the authors of \lstinline{opt_einsum} have given much more
thought to their heuristics, but it is not immediately available for all backends, and the
fact that my heuristics live above the backend and need to know nothing about the hardware
to at least make a *good* guess (maybe not the best) still makes it useful.  That said, 
on further study, we might find that \lstinline{opt_einsum} is just as easily put on top of other backends
as is my own code (eg, it already works on top of numpy and is default in pytorch).
If that is the case, and we find it is a hands-down winner over my own heuristic, then we can
just have evaluate() pass *all* of the arguments directly to \lstinline{opt_einsum}.  This might also
be left as a user choice.

Update:
\lstinline{opt_einsum} can now be enabled with the tensorly backend by the use of the following two lines, probably at the top-level:

  \lstinline{tensorly_plugins.use_opt_einsum()}

  \lstinline{qode.math.tensornet.backend_contract_path(True)}

See also:

  \url{https://tensorly.org/stable/modules/generated/tensorly.plugins.use_opt_einsum.html#tensorly.plugins.use_opt_einsum}

  \url{https://optimized-einsum.readthedocs.io/en/stable/install.html#conda}

I'm not a big fan of how this works because, IMO, all decisions about how to contract tensors should
be as local as possible.  Say, for example, you wanted to merge two codes that had been using different
backends for the tensors.  It will be a lot more convenient if you don't have to force them to agree on one.
This is why my tensors carry knowledge of their backends with them, and the only thing that matters when
two tensors meet is that they are the same (which is presumably the case within any given code domain).
Now there is this wart, where you insist that all tensors (even those with different backends) are either
going to use the frontend or the backend contraction path.  Probably this is ok for well over 90% of use
cases, but I would like to clean it up.  At the same time though, making the \lstinline{backend_contract_path}
option a feature of a backend is also not really (ie, could carry this option around with the tensors)
because it is not a feature of the backend, it is a feature of a contraction within the context of a 
backend.  We certainly do not want to insist that a tensor created with one way of determining the contraction
path cannot be contracted with a tensor with a different "opinion" on the way of determining the contraction
path.  At the moment then, this is the least of evils.  At least a user can dynamically switch back and forth.


\end{document}
